\section{Cartan bounds for quasisimple groups}\label{sec:families}

We now reduce the quasisimple Cartan problem to the odd-prime classical
calculations of the next three sections.  Throughout this section, the
coefficient prime is \(p\), the defining characteristic of a finite
reductive group is \(r\), and \(q\) is its defining field parameter.
For a positive integer \(a\), write \(a_p\) for its \(p\)-part.  For an
ordinary irreducible character \(\chi\) of a finite group \(H\), put
\[
 \operatorname{def}_p(\chi)=\frac{|H|_p}{\chi(1)_p}.
\]
This is the degree defect, expressed as an order rather than an exponent.
If \(\chi\) belongs to a block with defect group \(D\), then
\(\operatorname{def}_p(\chi)\leq |D|\).

\begin{theorem}[Quasisimple Cartan bound]\label{thm:quasisimple-cartan}
For every prime \(p\) and positive integer \(M\), there is a constant
\(C_{\mathrm{qs}}(p,M)\) such that every Cartan entry of every block of
\(kS\), where \(S\) is quasisimple and its defect groups have order at
most \(M\), is at most \(C_{\mathrm{qs}}(p,M)\).
\end{theorem}

The proof uses Theorem~\ref{thm:geometric-cartan} for bounded root
data.  For unbounded rank its inputs are Proposition~\ref{prop:runner-reduction}
and Proposition~\ref{prop:triangle-cartan}, with the character and block
identifications proved in Propositions~\ref{prop:label-degrees},
\ref{prop:single-cycle}, and~\ref{prop:core-blocks}.
Those propositions concern explicit reductive groups and do not use
Theorem~\ref{thm:quasisimple-cartan}.

\subsection{Restriction and reductive-group conventions}

We record the consequence of restriction that will be used when a
covering block can have large defect.

\begin{corollary}\label{fam:restriction}
Let \(N\lhd H\) with cyclic quotient, and let \(b\) be a block of
\(N\).  Suppose that every block of \(H\) covering \(b\) has one
simple module and all its decomposition numbers are one.
Then every decomposition number of \(b\) is at most one.
No defect bound on the covering blocks is required.  If the
original block \(b\) has bounded defect, its Cartan entries are
therefore bounded.
\end{corollary}

\begin{proof}
Apply Lemma~\ref{lem:restriction-rows} with \(L=c=u=1\).
The Cartan assertion uses only the Brauer--Feit bound for the
number of ordinary rows of \(b\).
\end{proof}

Write \(J=\mathbf J^F\), where \(\mathbf J\) is a connected reductive
group over \(\overline{\F}_r\), and write \(\mathbf J^*\) for its dual
group with its dual endomorphism.  We suppress the star on the latter
endomorphism.  If \(s\in\mathbf J^{*F}\) is semisimple, then
\(\mathcal E(J,s)\) denotes its ordinary Lusztig series.  When \(s\)
has order prime to \(p\), set
\[
 \mathcal E_p(J,s)
  =\bigcup_t\mathcal E(J,st),
\]
where \(t\) runs over the \(p\)-elements of
\(C_{\mathbf J^*}(s)^F\), with the usual conjugacy identifications.
Here and in the reductive reductions below \(r\ne p\); defining
characteristic is treated at the end of the section.  The
Brou\'e--Michel theorem states that this is a union of blocks
\cite{BroueMichel1989,CEBook}.

We use the following forms of standard reductive-group character
theory.  They also specify the hypotheses at every later application.
If \(Z(\mathbf J)\) and \(C_{\mathbf J^*}(s)\) are connected, Jordan
decomposition identifies \(\mathcal E(J,s)\) with the unipotent
characters of \(C_{\mathbf J^*}(s)^F\).  Its degree formula is
\begin{equation}\label{fam:jordan-degree}
 \chi(1)=
  |\mathbf J^{*F}:C_{\mathbf J^*}(s)^F|_{r'}\,\psi(1),
 \qquad
 \operatorname{def}_p(\chi)
  =\frac{|C_{\mathbf J^*}(s)^F|_p}{\psi(1)_p}.
\end{equation}
The torus-pairing rule identifies the pairings with corresponding
unipotent Deligne--Lusztig characters, multiplied by
\(\varepsilon_{\mathbf J}\varepsilon_{C_{\mathbf J^*}(s)}\), where
\(\varepsilon_{\mathbf H}=(-1)^{\rk_{\F_q}(\mathbf H)}\).
We use these two rules from \cite{LusztigCharacters}; see also
\cite[Equation~(3.1) and Theorem~7.1]{DigneMichel1990}.
A compatibility
statement for arbitrary nonuniform Levi induction is neither part of
these rules nor assumed here.  The particular induction matrices needed
below are established in Proposition~\ref{prop:single-cycle}.

Unipotent character degrees and unipotent induction calculations are
unchanged by central isogenies, with the standard identifications.
Connected isogenous reductive groups have the same rational order.
For a product of factors permuted by \(F\), the calculation on one
factor uses the composite endomorphism around its orbit.
A \emph{packet} is a Frobenius orbit of classical factors of a
connected centralizer, calculated on one factor with this composite
endomorphism.

For the ordinary classical Frobenius structures used below, we
make the unipotent identifications through connected kernels.
Choose an $F$-equivariant regular embedding
$\mathbf H\hookrightarrow\widetilde{\mathbf H}$ with connected
center and the same derived group.  Put
$H=\mathbf H^F$, $\widetilde H=\widetilde{\mathbf H}^F$, and
$A=\widetilde H/H$, an abelian group.  Twisting an upstairs
unipotent character $\widetilde\chi$ by a nontrivial linear
character of $A$ changes its dual parameter from $1$ to a
nontrivial central element.  Series disjointness and Clifford
theory therefore give
\[
 \langle\Res_H\widetilde\chi,\Res_H\widetilde\chi\rangle
 =\sum_{\nu\in\Irr(A)}
       \langle\widetilde\chi,\widetilde\chi\nu\rangle=1.
\]
Restriction is thus irreducible.  Every unipotent character of
$H$ occurs in some $R_T^{\mathbf H}(1)$; the torus restriction
formula lifts it to an upstairs unipotent character.  Two
upstairs unipotent characters with the same restriction differ
by a quotient twist, so series disjointness makes them equal.
This gives a bijection preserving degrees.  Next,
$\widetilde{\mathbf H}\to\mathbf H_{\mathrm{ad}}$ has connected
central kernel.  Lang's theorem makes it surjective on rational
points, and unipotent characters are trivial on that kernel.
These two morphisms compare the forms through their common
adjoint group; the relevant character and Levi diagrams are
compatible by \cite[Proposition~2.5(i), Theorem~9.1 and
Corollary~9.2]{DigneMichel1990}, with corresponding Levi preimages.
Both identifications are bijections, so individual split type-$D$
labels remain distinct.  Rationally permuted factors again use
the composite Frobenius on one factor.  This argument does not
apply the connected-kernel statement to a finite disconnected
central kernel; exceptional diagram isogenies remain in the
separate bounded-root-data case.

Restriction through a regular embedding, and restriction to a connected derived
subgroup, sends a Lusztig parameter to its dual projection.  Geometric
series suffice for these restriction statements; all later Jordan
matrix calculations are made in groups with connected centers and
connected semisimple centralizers.

For a good prime \(p\) and connected \(Z(\mathbf J)\), the integral
basic-set theorem \cite{GeckHiss} says that
\(\mathcal E(J,s)\) is an ordinary basic set for
\(\mathcal E_p(J,s)\).  Thus its restrictions to \(p\)-regular elements
form an integral basis of the Brauer character lattice of that block
union.  Its intersection with each block is nonempty and is a basic
set for that block.  We use this theorem for all primes in type~A,
and for odd primes in types~B, C, and D.  We do not apply it to type~B,
C, or D at \(p=2\).

\begin{lemma}[Central quotients in the Levi equivalence]
\label{fam:br-central}
Let \(s\in\mathbf J^{*F}\) be a semisimple \(p'\)-element, and let
\(\mathbf L^*\) be an \(F\)-stable Levi subgroup containing the full
algebraic centralizer \(C_{\mathbf J^*}(s)\).  The
Bonnaf\'e--Rouquier equivalence matches the blocks in
\(\mathcal E_p(J,s)\) with those in the corresponding summand of
\(L=\mathbf L^F\), preserving their Cartan matrices.
If \(Z\leq Z(\mathbf J)^F\) is contained in \(L\), this equivalence restricts
to the categories on which \(Z\) acts trivially.  In particular, it
gives the corresponding equivalences after taking a common central
quotient, including a quotient by a central \(p\)-subgroup.
\end{lemma}

\begin{proof}
The full-centralizer containment allows us to apply the integral
theorem \cite[Theorem~B$'$, Section~11.4]{BR}.
Its Deligne--Lusztig cohomology bimodule gives the asserted Morita
equivalence.  On the defining variety, left multiplication by a common
rational central element agrees with right multiplication by that
element.  Hence \(zm=mz\) on the equivalence bimodule, for every
\(z\in Z\).  The equivalence and its inverse consequently carry the
full subcategory annihilated by all \(z-1\) to the corresponding full
subcategory on the other side.  These are exactly the module
categories for the quotient algebras.  This argument uses equality of
the central actions; it does not use exactness of ordinary coinvariants
on arbitrary modules.  It applies over a splitting modular system and
over \(k\).  When \(Z\) has a \(p'\)-part, only blocks in its trivial
central-character sector survive this quotient.  The Cartan matrices
of surviving matched blocks agree under Morita equivalence, and
their largest elementary divisors give the equality of defect orders.
\end{proof}

\subsection{Linear and unitary groups}

Use the notation
\(\GL_n^+(q)=\GL_n(q)\), \(\GL_n^-(q)=\GU_n(q)\), and similarly
\(\SL_n^+(q)=\SL_n(q)\), \(\SL_n^-(q)=\SU_n(q)\).
A block in \(\mathcal E_p(J,1)\) will be called a unipotent block;
its ordinary characters need not all be unipotent.

\begin{proposition}\label{fam:type-a}
Fix \(p\) and \(M\).  To bound Cartan entries for blocks of defect
order at most \(M\) of all central quotients of
\(\SL_n^\epsilon(q)\), with \(r\ne p\), it suffices to bound Cartan
entries for unipotent blocks of \(\GL_m^\eta(Q)\) of bounded defect
order, where the new bound depends only on \(p,M\).
At \(p=2\), these latter blocks have bounded rank.
\end{proposition}

\begin{proof}
Let \(b\) be the original block and put
\[
 b_q=(q-\epsilon)_p,\qquad
 h=(\gcd(n,q-\epsilon))_p.
\]
We use \(b_q\) to distinguish this integer from the block \(b\).
Lift \(b\) to \(S=\SL_n^\epsilon(q)\), and cover the lifted block
in \(J=\GL_n^\epsilon(q)\).  A central \(p'\)-kernel merely selects
an averaging summand; a central \(p\)-kernel increases the defect
order by its order, at most \(h\).  Since \(J/S\) is cyclic of
order \(q-\epsilon\), every covering block \(B\) satisfies
\begin{equation}\label{fam:a-cover-defect}
 |D_B|\leq Mhb_q\leq Mb_q^2.
\end{equation}
We first allow \(b_q\) to be unbounded.

Let \(s\) be the \(p'\)-parameter of \(B\).  Its algebraic centralizer
is a rational Levi with fixed points
\begin{equation}\label{fam:a-centralizer}
 C_{\mathbf J^*}(s)^F
   \cong\prod_i\GL_{m_i}^{\epsilon^{d_i}}(q^{d_i}),
 \qquad n=\sum_i d_i m_i.
\end{equation}
Here \(d_i\) is the length of the relevant eigenvalue orbit.  The
basic-set theorem supplies a character of \(B\cap\mathcal E(J,s)\).
If its Jordan labels are partitions \(\lambda_i\vdash m_i\), the
partition degree formula and \eqref{fam:jordan-degree} give
\begin{equation}\label{fam:a-hook-defect}
 \operatorname{def}_p(\chi)
  =\prod_i\prod_{u\in\mathsf H(\lambda_i)}
       (q^{d_i u}-\epsilon^{d_i u})_p,
\end{equation}
where \(\mathsf H(\lambda_i)\) is the multiset of hook lengths, one
for each box.

If \(b_q>1\), lifting the exponent, applied to
\(x=\epsilon q\), yields
\begin{equation}\label{fam:a-lte}
 (x^j-1)_p\geq (x-1)_p(j)_p=b_q(j)_p.
\end{equation}
For odd \(p\) this is the usual equality.  For \(p=2\) and odd \(j\)
it is also an equality.  For even \(j\), the formula
\[
 v_2(x^j-1)=v_2(x-1)+v_2(x+1)+v_2(j)-1
\]
implies the inequality, including the case \(v_2(x-1)=1\).
Absolute values are understood when \(x<0\).
Consequently, if \(t=\sum_i m_i\), then
\begin{equation}\label{fam:a-multiplicity}
 b_q^t\prod_i(d_i)_p^{m_i}
 \leq\operatorname{def}_p(\chi)
 \leq Mhb_q\leq Mb_q^2.
\end{equation}
For \(b_q>2M\), this forces \(t\leq2\).  If some \(m_i=2\), it
is the only factor and \(n=2d_i\).  Hence
\(h\leq(2d_i)_p\leq2(d_i)_p\), so the same inequality gives
\[
 b_q^2(d_i)_p^2\leq2M(d_i)_p b_q,
 \qquad b_q(d_i)_p\leq2M,
\]
a contradiction.  Thus every \(m_i=1\): the centralizer is a torus.

Lemma~\ref{fam:br-central} now identifies \(B\) with a block of a
finite torus.  Such a block has one simple module, and all its
decomposition numbers are one.  The integral equivalence gives the
same decomposition matrix, up to labels, for \(B\).  This argument
applies to every covering block, so Corollary~\ref{fam:restriction}
bounds all decomposition numbers of the lifted block of \(S\) by
one.  Passing to the original central quotient only selects inflated
ordinary rows and the corresponding simple modules: the central
\(p\)-kernel acts trivially on every simple module.  The original
block still has defect order at most \(M\), so its number of ordinary
rows is bounded by the Brauer--Feit bound.  Its Cartan entries, which
are sums of products of entries in those rows, are bounded.  No
bound on the row count of \(B\), or on its defect in this case, has
been used.

We may therefore assume \(b_q\leq2M\).  Then
\eqref{fam:a-cover-defect} bounds \(|D_B|\) by \(4M^3\).
Apply Lemma~\ref{fam:br-central} to
\eqref{fam:a-centralizer}.  In this Levi the parameter \(s\) is
central.  Tensoring with its associated linear character identifies
its \(s\)-summand with its unipotent summand.  A block of the product
is a tensor product of blocks, and its defect order is the product
of their defect orders.  At most \(\log_p(4M^3)\) factors can have
positive defect.  Factors of defect zero have Cartan matrix \((1)\)
and have no effect on a Cartan bound.  This proves the stated
reduction, followed by cyclic restriction and central descent.

Finally suppose \(p=2\).  Here every defining parameter \(Q\) is
odd, and every hook factor \((Q^u-\eta^u)_2\) is at least two.
Every basic row of a unipotent block of \(\GL_m^\eta(Q)\) therefore
has degree defect at least \(2^m\).  Bounded block defect bounds
\(m\), and Theorem~\ref{thm:geometric-cartan} applies.
\end{proof}

\subsection{Regular classical root data}

The linear and unitary reduction has isolated unipotent blocks.
For the other classical families we must retain both a central
quotient and actual product decompositions of the Levis used
later.  The next construction supplies those groups before any
character or runner calculation is made.

The regular embedding must supply three kinds of structure.
The connectedness assertions are used in Jordan decomposition and,
at odd \(p\), in the basic-set theorem.  Integral Levi splittings
identify the actual product group algebras used in Harish--Chandra
induction.  Control of the residual central tori then keeps track of
their contributions to
degree defects, including after repeated extractions.
We begin with the ordinary irreducible root systems, in particular
\(D_n\) with \(n\geq4\); smaller ranks already fall under
Theorem~\ref{thm:geometric-cartan}.  Rank-one orthogonal tori
are nevertheless retained as residual factors and centralizer packets.

\begin{lemma}\label{fam:regular-datum}
For simply connected groups of types~B, C, and D there are regular
embeddings \(\mathbf S=[\mathbf J,\mathbf J]\leq\mathbf J\) with
central torus rank at most three with the following properties.
\begin{enumerate}
\item The groups \(\mathbf J\), its dual, and all their Levi
subgroups have connected centers and simply connected derived groups.
\item Every extraction of linear blocks on split hyperbolic
coordinates gives an actual direct product of these general linear
groups with a residual reductive group.  This holds for nested
extractions and in graph-twisted type~D.
\item Every residual factor \(\mathbf R\) has an inherited central torus
\(\mathbf T_R\), the image of the original connected center under
the Levi product projection, of rank at most three.  Its full
connected center has rank at most four.  The two tori agree except
for the rootless orthogonal datum \(D_1\), whose additional torus
direction is retained as an \(\mathrm{SO}_2^\pm\) packet.
With bounded residual semisimple rank, only finitely many residual
root data and normalized Frobenius actions occur.
\end{enumerate}
\end{lemma}

\begin{proof}
We first construct the ambient lattice and its Frobenius action.
We then split off the extracted general linear factors integrally.
The remaining calculations identify the inherited central torus
and show that bounded residual rank gives only finitely many
residual root data.

Let \(V=\Q^n\) with orthogonal coordinate vectors \(e_i\).  Denote
the coroot lattice and coweight lattice of the chosen classical root
system by \(Q^\vee\subset P^\vee\).  Explicitly,
\[
\begin{array}{c|c|c}
 &Q^\vee&P^\vee\\ \hline
 B_n&\{a\in\Z^n:\sum a_i\text{ is even}\}&\Z^n\\
 C_n&\Z^n&\Z^n\cup((\tfrac12,\ldots,\tfrac12)+\Z^n)\\
 D_n&\{a\in\Z^n:\sum a_i\text{ is even}\}
     &\Z^n\cup((\tfrac12,\ldots,\tfrac12)+\Z^n).
\end{array}
\]
Introduce a central space \(\Q^h\), with \(h=1,1,3\), respectively,
and define an injective homomorphism
\(\gamma:P^\vee/Q^\vee\to\Q^h/\Z^h\) by the following table.
Put \(\omega=(\tfrac12,\ldots,\tfrac12)\).
\begin{equation}\label{fam:lattice-glue}
\begin{array}{c|c|c}
 &\gamma(e_i)&\gamma(\omega)\\ \hline
 B_n&\tfrac12&\text{not needed}\\
 C_n&0&\tfrac12\\
 D_n,\ n\text{ even}&(\tfrac12,\tfrac12,0)
                         &(\tfrac12,0,\tfrac12)\\
 D_n,\ n\text{ odd}&(\tfrac12,\tfrac12,0)
                         &(\tfrac14,-\tfrac14,\tfrac12).
\end{array}
\end{equation}
All central entries are read modulo \(\Z^h\).  In odd rank~D,
\(2\gamma(\omega)=\gamma(e_i)\); in even rank the two displayed
classes of order two are independent.  Thus these prescriptions
respect exactly the relations in \(P^\vee/Q^\vee\).
Set
\begin{equation}\label{fam:Y}
 Y=\{(a,z)\in P^\vee\oplus\Q^h:
                   z+\Z^h=\gamma(a+Q^\vee)\},
 \qquad X=\Hom(Y,\Z).
\end{equation}
Take the usual roots on \(a\), extended by zero on \(z\), and the
usual coroots with zero central coordinate.  They define a root datum
on \((X,Y)\).  Indeed roots take integral values on \(P^\vee\),
coroots belong to \(Y\), and the classical reflection formulas
preserve the congruence because they change \(a\) by a coroot.

Injectivity of \(\gamma\) gives
\(Y\cap(V\oplus0)=Q^\vee\).  Since \(Y\to P^\vee\) is onto,
\(X\cap V^*=Q\), the root lattice.  These are the two saturation
conditions for simply connected derived group and connected center;
they remain true on dualizing.  For a Levi subsystem, its simple
roots form a subset of a simple-root basis, and likewise for
coroots.  Intersecting with their rational spans preserves these
saturation conditions.  This proves the assertions about every
Levi and its dual.  In particular their semisimple centralizers are
connected, by the connected-centralizer theorem for simply connected
derived groups.

In type~D the nontrivial diagram automorphism changes the sign of
the last coordinate of \(a\).  Extend it on \(z\) by exchanging
\(z_1,z_2\).  For even \(n\), it sends
\(\gamma(\omega)\) to \(\gamma(\omega)-\gamma(e_n)\);
the same calculation holds for odd \(n\).  It fixes
\(\gamma(e_i)\).  Thus it preserves \(Y\), and defines the
graph-twisted Frobenius on this datum.  The action on the central
space, after the scalar \(q\) is removed, has order at most two.

\smallskip\noindent\emph{Integral Levi products.}
For the splitting assertion, orient the extracted hyperbolic
coordinates by signs \(\sigma_i\in\{1,-1\}\).  Let \(c\) be the
entry \(\gamma(e_i)\) in the table, with the displayed representative,
and set
\[
 f_i=(\sigma_i e_i,c),\qquad
 x_i(a,z)=\sigma_i a_i+\ell(z),\qquad
 \ell(z)=
 \begin{cases}
  0&\text{in type B},\\
  z_1&\text{in type C},\\
  z_3&\text{in type D}.
 \end{cases}
\]
Then \(f_i\in Y\), and \(x_i\in X\): the possible half-integral
part of \(a_i\) is canceled by the indicated central coordinate.
Moreover \(\ell(c)=0\), so
\begin{equation}\label{fam:split-lattice}
 x_i(f_j)=\delta_{ij},\qquad
 Y=\bigoplus_{i\in I}\Z f_i\ \oplus\
               \bigcap_{i\in I}\ker x_i.
\end{equation}
Within an extracted block, its roots are \(x_i-x_j\) and its
coroots are \(f_i-f_j\).  They are exactly the root datum of the
corresponding general linear group.  The residual roots and
coroots lie on the complementary summand.  Thus
\eqref{fam:split-lattice} is a product decomposition of root data,
not merely of rational root spaces.

For a rational split hyperbolic extraction, conjugate its split
cocharacters into standard position on a maximal split torus.
The construction is then \(F\)-compatible.  In nonsplit type~D,
the extracted directions are in split hyperbolic planes; the
remaining diagram involution acts on the residual coordinates and
on \(z\) as above.  The same construction applies after each
extraction, proving the nested assertion.

\smallskip\noindent\emph{Residual center and root data.}
The general linear factors have now been split off.  We retain
the image of the original center separately from any extra torus
direction in the residual group.
To identify this inherited center, let \(m=|I|\).  The projection
of an original central vector \((0,z)\) to the residual summand is
\((0,z)-\sum_{i\in I}x_i(0,z)f_i\).  Its remaining
\(a\)-coordinates are zero, its discarded coordinates are
\(a_i=-\sigma_i\ell(z)\), and its central coordinate is
\[
 z'=(\id-mc\ell)z.
\]
Since \(\ell(c)=0\), the central map has inverse \(\id+mc\ell\)
and determinant one.  It is \(F\)-equivariant, so its image
\(\mathbf T_R\) is a torus of rank \(h\) with the same rational
Frobenius representation as the original center.  The central
lattices are commensurable with the coordinate lattice only at
the prime two, by the graph congruences.  Together with the
determinant-one map this shows that the isogeny onto
\(\mathbf T_R\) has kernel of 2-power order.  For odd \(p\)
it therefore identifies the Sylow \(p\)-subgroups on rational
points.  The original center may also project to the general
linear centers; it need not lie solely in the residual factor.

If the residual root system spans the remaining coordinate
space, then \(\mathbf T_R=Z^\circ(\mathbf R)\).  In the
rootless datum \(D_1\), the remaining coordinate instead adds
one central torus dimension with Frobenius eigenvalue \(q\)
or \(-q\).  Thus the full residual central rank is at most
\(h+1\leq4\).  This extra torus is the rank-one orthogonal
packet, separate from the inherited central contribution.

Finally, on the residual summand the equations \(x_i=0\) express
every discarded coordinate as \(a_i=-\sigma_i\ell(z)\).
Substitution into \eqref{fam:Y} leaves congruences with uniformly
bounded denominators, dividing four.  Independently of how many
coordinates were discarded, the resulting lattice lies in a
fixed small-denominator coordinate lattice and contains a fixed
power-of-two multiple of the integral coordinate lattice.  In
bounded remaining dimension there are only finitely many intermediate
lattices.  The residual signed-coordinate and diagram actions are
also drawn from a finite set in that dimension.  This proves the
last assertion, including the rank-one orthogonal residual tori.
\end{proof}

\subsection{Reduction to the two sign eigenspaces}

For an original regular ambient group of
Lemma~\ref{fam:regular-datum}, before any split extraction, put
\[
 T_0=Z^\circ(\mathbf J)^F,\qquad Z_p=O_p(T_0).
\]
Here \(O_p(T_0)\) denotes its unique Sylow \(p\)-subgroup.  A block
of \(J/Z_p\) with defect order at most \(M\) lifts to a block of
\(J\) with defect order at most \(M|Z_p|\).  We call this a relative
defect bound.  It will always be accompanied by the requirement
that the final Cartan bound is for the block of \(J/Z_p\).
For odd \(p\), the basic rows at a semisimple \(p'\)-parameter
\(s\) descend to this quotient.  Indeed each such row occurs in
a Deligne--Lusztig character \(R_T^{\mathbf J}(\theta)\) for a torus
character \(\theta\) dual to \(s\).  Torus duality gives
\(|\theta|=|s|\), and \(Z(\mathbf J)^F\) acts on every constituent
by the restriction of \(\theta\) to \(Z(\mathbf J)^F\)
\cite[Corollary~1.22 and Section~5.21]{DL}.
Their central characters therefore have \(p'\)-order and are trivial
on \(Z_p\).  Simple modules also factor through \(J/Z_p\), by
Lemma~\ref{lem:central-transfer}.  Thus the descended rows form the
same integral basic set on the corresponding quotient block union,
and their degree defects there are
\(\operatorname{def}_p(\chi)/|Z_p|\).

For a later residual factor, the central contribution in its
packet degree formula will be \(\mathbf T_R\), rather than its
possibly larger full connected center.

We also fix the spectral convention in the following reduction.
Project a semisimple \(p'\)-parameter \(s\) to the adjoint dual
group and represent it in natural symplectic or orthogonal
eigenvalue coordinates.  The passage to a natural isometry group
requires at most a central 2-isogeny.  We call this eigenvalue
multiset the \emph{normalized natural spectrum}; its ambiguity,
and the ambiguity in its rational Frobenius action, is at most a
common sign.  In particular, containment in \(\{1,-1\}\) is
independent of this choice.

\begin{proposition}\label{fam:central-classical}
To prove Theorem~\ref{thm:quasisimple-cartan} for the
cross-characteristic families of types~B, C, and D, it suffices to
prove the following assertion, together with the type~A bounds:
for the groups of Lemma~\ref{fam:regular-datum}, blocks of
\(J/Z_p\) of bounded defect have bounded Cartan entries whenever
their \(p'\)-parameter has normalized natural spectrum contained
in \(\{1,-1\}\).
\end{proposition}

\begin{proof}
The centers on points of the simply connected covers of types~B, C, and~D have
order at most four.  By Lemma~\ref{lem:central-transfer}, it
suffices to treat these covers and their quotients by central
\(p\)-subgroups, with defect bounds changed by a factor at most four.
Embed \(S=\mathbf S^F\) into \(J\) as in
Lemma~\ref{fam:regular-datum}, and put \(K=S\cap Z_p\).
The central isogeny
\(\mathbf S\times Z^\circ(\mathbf J)\to\mathbf J\), followed
by Lang's theorem, gives
\[
 [J:S T_0]\leq4,\qquad
 [J/Z_p:S/K]_p\leq4.
\]
Thus every block of \(J/Z_p\) covering a given block of \(S/K\)
has bounded defect.  A simple module of \(J/Z_p\) restricts to
\(S/K\) with bounded multiplicity: the central group \(T_0/Z_p\)
acts by scalars on an inertia type, and the remaining inertia
quotient has order at most four.  Lemma~\ref{lem:restriction-rows}
therefore transfers a Cartan bound for these covering blocks
down to \(S/K\).  Central transfer then recovers \(S\) and its
required central quotients.

Consider one such covering block, with \(p'\)-parameter \(s\).
Eigenvalues different from \(1,-1\) occur in
reciprocal pairs and give linear root subsystems in the connected
centralizer.

There is an \(F\)-stable Levi \(\mathbf M^*\) containing this full
centralizer and retaining the full ambient classical subsystem on
the \(\{1,-1\}\)-coordinates.  To construct it, take independent
torus directions with a scalar on each non-sign eigenvalue space,
its inverse on the reciprocal space, and zero direction on the
remaining coordinates, and centralize this torus.  The span of
these directions is stable under Frobenius, which permutes the
pairs and may invert them or multiply all normalized eigenvalues
by a common sign.  The full centralizer is contained because its
roots do not join different reciprocal packets.  The centralizer
is connected by Lemma~\ref{fam:regular-datum}.

Apply Lemma~\ref{fam:br-central}, also modulo \(Z_p\), and let
\(\mathbf M\) be the corresponding primal Levi.  Its connected
derived subgroup has simply connected factors.  The normal subgroup
\[
 N=[\mathbf M,\mathbf M]^F/
       ([\mathbf M,\mathbf M]^F\cap Z_p)
       \ \lhd\ M/Z_p
\]
has abelian quotient.  The intersection removed here lies in
\(S\), so has order at most four.  Restrict a block of \(M/Z_p\)
to \(N\), and lift across that bounded central intersection.
The covered blocks on \([\mathbf M,\mathbf M]^F\) have bounded
defect.  This group is a direct product of simply connected
linear or unitary groups and possibly one residual group of type~B, C, or~D.
Only boundedly many of its block factors can have positive defect.

On the residual factor, dual projection keeps precisely the roots
on the sign coordinates.  Projection commutes with taking the
prime-to-\(p\) part of a semisimple element.  The residual block
therefore has a geometric \(p'\)-parameter with normalized spectrum
in \(\{1,-1\}\).  Embed this factor regularly again as above;
all lifted parameters retain this adjoint-spectrum property.
The stated special assertion bounds its blocks.  The type~A
bounds handle the other positive-defect factors.  Tensor products,
central transfer, and Lemma~\ref{lem:abelian-transfer} then bound
the blocks of \(M/Z_p\), and hence the original block.
If no non-sign eigenvalue is present, no Levi reduction is made.
Otherwise the residual classical rank is strictly smaller.  In
either case this argument invokes only the stated special
assertion, not the general quasisimple conclusion.
\end{proof}

\subsection{The prime two}

\begin{proposition}\label{fam:two}
At \(p=2\), the blocks in the special assertion of
Proposition~\ref{fam:central-classical} have bounded ambient
semisimple rank.  They consequently have bounded Cartan entries.
\end{proposition}

\begin{proof}
Here \(q\) is odd.  The adjoint projection of the odd-order element
\(s\) has a natural lift with spectrum in \(\{1,-1\}\).  Choose
the lift of odd order.  Its spectrum is then \(\{1\}\), so \(s\)
is central.  Tensoring by the corresponding central linear
character removes \(s\) from the calculation.

Let \(B\) be the lifted block of \(J\), whose defect order is at
most \(M|Z_2|\), and choose any \(\chi\in\Irr(B)\).  Its
semisimple parameter is \(st\), where \(t\) is a 2-element.
Put \(\mathbf C=C_{\mathbf J^*}(t)\), which is connected.
The degree formula identifies \(\operatorname{def}_2(\chi)\)
with the degree defect of its unipotent correspondent \(\psi\)
in \(\mathbf C^F\).  Unipotent characters are trivial on the
center: they occur in Deligne--Lusztig characters induced from
trivial torus characters, on which that center acts trivially.
In particular the central subgroup
\[
 A=\langle t,O_2(Z^\circ(\mathbf J^*)^F)\rangle
\]
lies in the kernel of \(\psi\), and hence
\[
 |A|\leq\operatorname{def}_2(\chi)\leq M|Z_2|.
\]
The two ambient central tori have the same rational Frobenius
eigenvalues under duality and isogeny, so
\(|O_2(Z^\circ(\mathbf J^*)^F)|=|Z_2|\).
It follows that the order of \(t\) modulo the ambient central
torus is at most \(M\).

Its adjoint image therefore has bounded order.  A lift to the
natural isometry group has order at most twice this bound.  Its
normalized eigenvalues belong to a set of roots of unity of
bounded order.  Both the number of distinct eigenvalues and
the lengths of their Frobenius orbits are consequently bounded;
the possible common-sign ambiguity only enlarges these bounds
by a factor of two.

Up to central isogeny, the rational root datum of \(\mathbf C\)
is the product of its ambient central torus and its eigenvalue
packets.  A reciprocal non-sign packet gives a linear or unitary
factor, and a sign packet gives a factor of type~B, C, or~D.  Factors permuted
by Frobenius use their composite field parameter.  Rank-one
orthogonal tori \(\mathrm{SO}_2^\pm\) are included as rank-one
type~D packets.  Thus rational orders and unipotent degrees give
\begin{equation}\label{fam:two-product}
 \frac{\operatorname{def}_2(\chi)}{|Z_2|}
   =\prod_i\frac{|H_i|_2}{\psi_i(1)_2}\leq M,
\end{equation}
where \(\psi_i\) is unipotent and \(H_i\) is the corresponding
packet group.  This equality is an order and degree calculation;
no direct-product assertion about the finite centralizer is needed.

A linear or unitary packet of dimension \(n\) contributes at least
\(2^n\), by the hook formula.  For completeness, the symbol formula
gives a similarly explicit estimate for every packet of type~B, C, or~D
and positive rank \(n\).  Write its two strictly increasing beta rows
as \(X,Y\), let \(u=|X|+|Y|\), and let \(c=|X\cap Y|\).
The rank formula is
\[
 n=\sum_{a\in X}a+\sum_{a\in Y}a
                  -\left\lfloor\frac{(u-1)^2}{4}\right\rfloor.
\]
Count all hooks and cohooks of positive length: these are pairs
of a bead and a strictly lower vacant position in its own row
or the opposite row.  Before vacancies are imposed there are
\(2\sum a\) possible pairs.  The number ending at occupied
positions is \(\binom u2-c\).  Therefore their number is
\begin{align*}
 H&=2\left(n+\left\lfloor\frac{(u-1)^2}{4}\right\rfloor\right)
                 -\binom u2+c\\
  &=2n-\lfloor u/2\rfloor+c.
\end{align*}
The unipotent symbol degree formula \cite[Section~3.3]{LMT}
expresses the 2-defect as the product of
\((q_i^j-1)_2\) over hooks and \((q_i^j+1)_2\) over cohooks,
times a power of two whose exponent is at least
\(\lfloor(u-1)/2\rfloor-c\).
For unequal rows this exponent is exactly the displayed one;
for equal rows in degenerate type~D the exponent is zero, which
is larger.  Every positive-length hook or cohook contributes
at least two, whence
\[
 v_2\!\left(\frac{|H_i|}{\psi_i(1)}\right)
 \geq 2n-\lfloor u/2\rfloor+\lfloor(u-1)/2\rfloor
 \geq 2n-1.
\]
This argument allows a negative displayed prefactor exponent;
it is the total exponent that is estimated.  The degenerate
case gives at least \(2n\).  For rank-one orthogonal tori the
defect is directly \((q_i-1)_2\) or \((q_i+1)_2\), and is at
least two.  Empty packets contribute one.

Equation~\eqref{fam:two-product} now bounds the effective rank
of every packet and the number of nonempty packets: each
positive-rank packet contributes at least \(2^{n_i}\), so
\(\sum_i n_i\leq\log_2 M\).  Their
Frobenius orbit lengths were already bounded, so it also bounds
the actual ambient semisimple rank.  The central rank is at
most three.  Lemma~\ref{fam:regular-datum} gives finitely many
root data, and Theorem~\ref{thm:geometric-cartan}, applied after
the central quotient, proves the result.
\end{proof}

\subsection{Odd-prime packets and the classical reduction}

The prime-two case is now reduced to fixed root data.  At an odd
prime the rank can remain unbounded; we instead express the
degree defect as a product of contributions from at most two
classical packets and a controlled central torus.  These are the
inputs for the label, runner and endpoint arguments that follow.

The next lemma describes the sign packets and their split Levi
factors, on which the subsequent label, runner, and endpoint
arguments operate.

\begin{lemma}\label{fam:packets}
Suppose that \(p\) is odd and the normalized spectrum of the
\(p'\)-parameter \(s\) is contained in \(\{1,-1\}\), in a
regular group of Lemma~\ref{fam:regular-datum}.
Apart from the ambient central torus, its connected dual
centralizer has at most two classical factors.  In natural dual
type~C they are of types C and C; in natural dual type~B they are
of types B and D; in natural dual type~D they are of types D
and D.  Empty factors are omitted and rank-one type~D is toral.
The factors have parameters \(q_i=q\), unless two like factors
are exchanged by Frobenius, in which case there is one packet
with parameter \(q_i=q^2\).

If split hyperbolic blocks of sizes \(b_j\) are extracted in
one packet with parameter \(q^\delta\), \(\delta\in\{1,2\}\),
the corresponding primal split Levi has actual factors
\(\GL_{\delta b_j}(q)\).  On each such factor the parameter
has a single eigenvalue orbit of degree \(\delta\) and
multiplicity \(b_j\).  The residual parameter retains the
same sign-packet description.  These assertions persist under
nested extractions with earlier factors held fixed.
\end{lemma}

\begin{proof}
Inspect the classical roots on the natural coordinate spaces.
In symplectic type, the roots evaluating trivially on \(s\)
are the type~C roots within its two sign eigenspaces.  In odd
orthogonal type, the positive eigenspace contains the distinguished
odd coordinate, giving a B factor, and the negative eigenspace
gives a D factor.  In even orthogonal type both eigenspaces
have even dimension, giving two D factors.  Their form signs
determine their split or graph-twisted rational structures and
must be retained.  In defining characteristic two the two signs
coincide, leaving a single packet.  Frobenius can exchange only
like factors; hence its orbit lengths here are at most two.

On a chosen split hyperbolic block, take a cocharacter that is
constant on its polarized half and on the Frobenius translates
of that half, has the opposite value on the reciprocal halves,
and is zero elsewhere.  Independent choices for the extracted
blocks define a split torus.  Its centralizer in the ambient
dual group is a split Levi, and its intersection with
\(C_{\mathbf J^*}(s)\) is exactly the desired split Levi of
the packet.  A packet coordinate has \(\delta\) ambient
coordinates in its Frobenius orbit, so the ambient linear
block has size \(\delta b_j\).

By Lemma~\ref{fam:regular-datum}, dualizing this construction
gives an actual product with \(\GL_{\delta b_j}(q)\) on the
primal side.  The centralizer roots on such a block form one
linear system of rank \(b_j-1\) on each of its \(\delta\)
Frobenius translates.  Equivalently its eigenvalues have one
orbit of degree \(\delta\) and multiplicity \(b_j\).
On the unextracted coordinates the original root evaluations
are unchanged, proving the residual assertion.  The independent
cocharacters can be chosen with all earlier extracted coordinates
fixed, which proves the assertion for a chain of extractions.
\end{proof}

The order and degree calculation used at \(p=2\) also gives,
for a basic row \(\chi\in\mathcal E(J,s)\),
\begin{equation}\label{fam:packet-defect}
 \frac{\operatorname{def}_p(\chi)}{|Z_p|}
       =\prod_i\frac{|H_i|_p}{\psi_i(1)_p}.
\end{equation}
Here the groups \(H_i\) are the packets of
Lemma~\ref{fam:packets}; their orthogonal signs are part of the
data.  This identity follows by decomposing the rational
reflection space into the packet spaces and the ambient central
space, then using central-isogeny invariance of orders and
unipotent degrees.  The partition and symbol formulas for its
right-hand side are made explicit in
Proposition~\ref{prop:label-degrees}.
For a residual factor \(\mathbf R\), the corresponding identity is
\[
 \operatorname{def}_p(\chi_R)
  =|\mathbf T_R^F|_p
        \prod_i\frac{|H_i|_p}{\psi_i(1)_p}.
\]
The product includes its toral \(D_1\) packet if present;
using the inherited torus ensures that this packet is counted once.

\begin{lemma}\label{fam:central-order}
For odd \(p\), the original central subgroup and every inherited
residual central torus satisfy
\[
 |Z_p|,\ |\mathbf T_R^F|_p
       \leq\bigl((q-1)_p(q+1)_p\bigr)^3.
\]
The full residual connected center satisfies the analogous bound
with exponent four.  Let \(q_i=q^\delta\), \(\delta\in\{1,2\}\),
be an original packet parameter,
\(e_i=\ord_p(q_i)\), and put
\[
 (d_i,\eta_i)=
 \begin{cases}
  (e_i,1)&e_i\text{ odd},\\
  (e_i/2,-1)&e_i\text{ even}.
 \end{cases}
 \qquad
 \alpha_i=(q_i^{d_i}-\eta_i)_p.
\]
Then \(|Z_p|\) and \(|\mathbf T_R^F|_p\) are at most
\(\alpha_i^3\), while \(|Z^\circ(\mathbf R)^F|_p\) is at most
\(\alpha_i^4\).  Thus a bound on \(\alpha_i\) bounds both
inherited and full residual central \(p\)-orders.
\end{lemma}

\begin{proof}
On the inherited central rational space Frobenius is \(q\theta\),
where \(\theta^2=1\), and the dimension is at most three.
The full residual center adds at most one coordinate with
eigenvalue \(q\) or \(-q\).  Its rational order is therefore
\((q-1)^a(q+1)^b\), with \(a+b\leq3\) for the inherited
torus and \(a+b\leq4\) for the full center.  These orders
are unchanged by isogeny of their integral lattices.
If either \(p\)-part is nontrivial, then \(p\mid q-1\) or
\(p\mid q+1\), and for odd \(p\) at most one occurs.
For \(\delta=1\), the relevant \(\alpha_i\) is exactly the
nontrivial one of these two \(p\)-parts.  For \(\delta=2\),
we have \(e_i=1\) and
\(\alpha_i=(q^2-1)_p=(q-1)_p(q+1)_p\).
If both \(p\)-parts are trivial the estimates are immediate;
otherwise the preceding calculation proves them.
\end{proof}

We now identify precisely how the later results finish this
section.  For a block of relative defect at most \(M\),
\eqref{fam:packet-defect} and
Proposition~\ref{prop:label-degrees} bound the total relevant
hook or cohook weight.  A positive weight bounds the corresponding
cyclotomic \(p\)-part \(\alpha_i\), so
Lemma~\ref{fam:central-order} bounds the ambient central
\(p\)-part whenever it is needed.  At weight zero, the core
arguments handle the central quotient directly.  The parameter
on each Levi is the specified rational parameter, rather than
a sum over its several possible preimages from the ambient
group.  Propositions~\ref{prop:single-cycle} and
\ref{prop:core-blocks} establish, respectively, the required
ordinary matrices on these labels and the fact that the
projectors used are projectors to unions of blocks.

Proposition~\ref{prop:runner-reduction} then transfers Cartan
bounds from two kinds of endpoints.  The first has bounded
active rank.  Its extracted factors are of order prime to
\(p\), split off by Lemma~\ref{fam:regular-datum}, and contribute
only defect-zero blocks.  Its residual root data form a finite
set by that lemma, so Theorem~\ref{thm:geometric-cartan} applies.
The second endpoint has bounded ordinary Harish--Chandra rank
above a possibly arbitrarily large cuspidal triangle.
Proposition~\ref{prop:triangle-cartan} gives the required bound
there.  The same two endpoint arguments cover unipotent
\(\GL\) and \(\GU\) blocks in
Proposition~\ref{fam:type-a}.  Thus these later propositions
prove both of the assertions left open by
Propositions~\ref{fam:type-a} and~\ref{fam:central-classical}
for every odd \(p\).

\subsection{Completion over all quasisimple families}

\begin{proof}[Proof of Theorem~\ref{thm:quasisimple-cartan}]
We use the classification of finite simple groups and the standard
Schur multiplier descriptions \cite{Wilson,MalleTesterman}.
Sporadic groups and the finitely
many exceptional multiplier stems form a bounded collection of
finite groups; all their blocks may be included in the bound.

For alternating groups, Scopes' theorem \cite{Scopes} gives
finitely many Morita classes of symmetric-group blocks of each
fixed weight.  Bounded defect order bounds this weight, since
the defect group is a Sylow \(p\)-subgroup of the symmetric
group on \(pw\) letters for a block of weight \(w\).
For the double covers, Kessar's finiteness theorem
\cite{KessarSymmetric} supplies the same bounded-defect
conclusion for spin blocks at odd primes; the nonspin blocks
factor through the symmetric group.  At \(p=2\), quotient by
the central involution and use central transfer.  Every block
of a double cover of an alternating group is covered in the
corresponding double cover of the symmetric group.  Its
covering defect increases by at most a factor two, and
restriction across this cyclic quotient has multiplicity one.
Lemma~\ref{lem:restriction-rows} therefore bounds its decomposition
numbers and Cartan entries.  The same argument without the
central cover treats the alternating group itself.

Consider groups of Lie type in defining characteristic
\(p=r\).  For their standard simply connected covers, the
defining-characteristic block theorem \cite{HumphreysDefining}
says that every positive-defect block has full defect.  The
standard central kernels on rational points have order prime
to \(r\) \cite[Proposition~3.8]{BMO}.
Consequently a positive-defect block on a central
quotient has defect equal in order to a Sylow \(r\)-subgroup
of the cover.  Bounded Sylow \(r\)-order bounds both its
defining field parameter and its rank, and hence its group
order.  The bounded list of small exceptions is already harmless.
\mbox{Defect-zero blocks have Cartan matrix \((1)\).}

For \(r\ne p\), every bounded-rank Lie-type family, including
all exceptional types, Suzuki and Ree groups, and triality,
is covered by Theorem~\ref{thm:geometric-cartan}.  Its hypotheses
allow the finitely many root data, diagram actions, and
exceptional diagram isogenies that occur, as well as their
central quotients.  Thus only the ordinary unbounded-rank
classical families remain.

For these, Proposition~\ref{fam:type-a} gives the type~A
reduction and handles the large central lift without imposing
a bounded defect upstairs.  The bounded-lift case at \(p=2\)
has bounded rank.  Proposition~\ref{fam:central-classical}
reduces types~B, C, and D to the regular sign-spectrum case,
which Proposition~\ref{fam:two} handles at \(p=2\).
For odd \(p\), Propositions~\ref{prop:runner-reduction} and
\ref{prop:triangle-cartan}, combined with the packet estimates above,
supply the remaining bounds for general linear, unitary, and regular
groups of types~B, C, and~D.  The restriction,
central-quotient, product, and abelian-extension arguments
already proved then return these bounds to every quasisimple
group in the corresponding families.  All changes in the
defect bound depend only on \(p,M\), completing the proof.
\end{proof}
